\structuralpart{ЗАКЛЮЧЕНИЕ}
В ходе работы получены следующие результаты:
\begin{itemize}
\item выбран метод Оцу, определены полутоновое исходное изображение и бинарная маска результата, зафиксированы правила обработки однотонного входа и одинаковых максимумов критерия;
\item разработаны последовательный и параллельный алгоритмы, схемы единственного последовательного потока, диспетчера и одного вида рабочих потоков с двумя зависимыми фазами;
\item выполнены реализации на языке Си11, проверены типовые и граничные случаи, совпадение результатов обеих версий и обработка неверного ввода;
\item выполнены замеры реального времени для пяти размеров изображений, самостоятельной последовательной версии и параллельной версии при 1, 2, 4, 8, 16 и 32 рабочих потоках;
\item установлено, что для небольших входов параллельная версия замедляет расчёт, а для $4096\times4096$ минимальное наблюдавшееся среднее время получено при восьми рабочих с ускорением 1,79; дальнейшее увеличение числа потоков устойчивого выигрыша не дало.
\end{itemize}
